\documentclass[a4paper,12pt]{article}

% ========================================================================
% Pacotes
% ========================================================================
\usepackage[margin=2.5cm]{geometry}
\usepackage[utf8]{inputenc}
\usepackage[T1]{fontenc}
\usepackage[brazil]{babel}
\usepackage{tikz}
\usepackage{amsmath}
\usepackage{amssymb}
\usepackage{graphicx}
\usepackage{float}
\usepackage{listings}
\usepackage{xcolor}
\usepackage{caption}
\usepackage{fancyhdr}

% TikZ
\usetikzlibrary{
    arrows.meta, positioning, calc, shapes.geometric,
    fit, decorations.markings, patterns
}

% ========================================================================
% Estilos TikZ compartilhados
% ========================================================================
\tikzset{
    block/.style = {
        rectangle, draw, very thick,
        minimum width=3.5cm, minimum height=1.8cm,
        font=\sffamily\small, align=center, text width=3cm
    },
    wide block/.style = {
        block, minimum width=5cm, text width=4.5cm
    },
    small block/.style = {
        rectangle, draw, thick,
        minimum width=2cm, minimum height=1cm,
        font=\sffamily\footnotesize, align=center
    },
    operator/.style = {
        circle, draw, very thick,
        minimum size=9mm, font=\sffamily\large,
        inner sep=0pt, fill=white
    },
    dot/.style = {
        circle, fill=black, minimum size=3pt, inner sep=0pt
    },
    port/.style = {font=\sffamily\bfseries\small},
    signal/.style = {font=\sffamily\footnotesize, fill=white, inner sep=1pt},
    bus/.style = {font=\sffamily\scriptsize},
    entity border/.style = {
        rectangle, draw, dashed, thick,
        rounded corners=3pt, inner sep=8pt
    },
    register/.style = {
        rectangle, draw, thick,
        minimum width=1.0cm, minimum height=0.8cm,
        font=\sffamily\footnotesize, align=center
    },
    data arrow/.style = {-Stealth, thick},
    ctrl arrow/.style = {-Stealth, thin, dashed},
    const/.style = {font=\sffamily\itshape\footnotesize},
    annotation/.style = {font=\sffamily\small, align=left},
    >=Stealth, thick, font=\sffamily
}

% ========================================================================
% Estilo VHDL
% ========================================================================
\definecolor{vhdlkeyword}{RGB}{0, 0, 160}
\definecolor{vhdlcomment}{RGB}{0, 128, 0}
\definecolor{vhdlbg}{RGB}{248, 248, 248}
\definecolor{vhdlnumber}{RGB}{120, 120, 120}

\lstdefinestyle{vhdlstyle}{
    language=VHDL,
    backgroundcolor=\color{vhdlbg},
    basicstyle=\ttfamily\scriptsize,
    keywordstyle=\color{vhdlkeyword}\bfseries,
    commentstyle=\color{vhdlcomment}\itshape,
    numberstyle=\tiny\color{vhdlnumber},
    numbers=left,
    numbersep=5pt,
    frame=single,
    rulecolor=\color{black!30},
    breaklines=true,
    tabsize=4,
    showstringspaces=false,
    captionpos=b,
    xleftmargin=2em,
    framexleftmargin=1.5em,
    morekeywords={data_t, data_long_t, data_array_t, signed, std_logic,
                  std_logic_vector, to_signed, shift_right, resize}
}

\lstdefinestyle{logstyle}{
    backgroundcolor=\color{black!5},
    basicstyle=\ttfamily\scriptsize,
    frame=single,
    rulecolor=\color{black!30},
    breaklines=true,
    captionpos=b,
    xleftmargin=1em,
    framexleftmargin=0.5em,
}

\definecolor{pykeyword}{RGB}{0, 100, 0}
\definecolor{pystring}{RGB}{160, 0, 0}
\definecolor{pycomment}{RGB}{100, 100, 100}
\definecolor{pybg}{RGB}{252, 252, 248}

\lstdefinestyle{pystyle}{
    language=Python,
    backgroundcolor=\color{pybg},
    basicstyle=\ttfamily\scriptsize,
    keywordstyle=\color{pykeyword}\bfseries,
    stringstyle=\color{pystring},
    commentstyle=\color{pycomment}\itshape,
    numberstyle=\tiny\color{vhdlnumber},
    numbers=left,
    numbersep=5pt,
    frame=single,
    rulecolor=\color{black!30},
    breaklines=true,
    tabsize=4,
    showstringspaces=false,
    captionpos=b,
    xleftmargin=2em,
    framexleftmargin=1.5em,
}

% ========================================================================
% Cabeçalho
% ========================================================================
\pagestyle{fancy}
\fancyhf{}
\fancyhead[L]{\sffamily\small Rede Neural DPD em VHDL}
\fancyhead[R]{\sffamily\small Atividade 13}
\fancyfoot[C]{\thepage}

% ========================================================================
% Título
% ========================================================================
\title{Relat\'orio -- Atividade 13 \\
       \large Integra\c{c}\~ao do Est\'agio~1 (\texttt{Top\_NeuralNet\_S1})}
\author{Jo\~ao Pedro Hinning}
\date{\today}

% ========================================================================
% Documento
% ========================================================================
\begin{document}

\maketitle
\thispagestyle{fancy}

% =======================================================================
\section{Vis\~ao Geral}
% =======================================================================

Esta atividade documenta a integra\c{c}\~ao e valida\c{c}\~ao do Est\'agio~1 completo da rede
neural DPD, implementado na entidade \texttt{Top\_NeuralNet\_S1}. O Est\'agio~1 combina todos
os blocos desenvolvidos nas atividades anteriores em um \'unico pipeline coerente:

\begin{itemize}
    \item \textbf{1$\times$ \texttt{Input\_Processing\_Block}:} Convers\~ao IQ$\,\rightarrow\,$polar
          via CORDIC (lat\^encia de 40~ciclos), desembrulhamento de fase e normaliza\c{c}\~ao.
    \item \textbf{2$\times$ \texttt{LUT\_Trig}:} Uma inst\^ancia para cosseno (\texttt{mode\_sin='0'})
          e outra para seno (\texttt{mode\_sin='1'}), operando sobre $\Delta\theta(n)$
          (\texttt{theta\_norm\_out}) com pipeline de 4~ciclos.
    \item \textbf{3$\times$ \texttt{Shift\_Register}:} Buffers deslizantes de $M+1 = 6$~taps
          para os canais de envelope, cosseno e seno. Os tr\^es registradores s\~ao habilitados
          sincronamente pelo sinal de escrita alinhado.
\end{itemize}

As sa\'idas do est\'agio s\~ao os tr\^es vetores de 6~taps ---
\texttt{buf\_env\_out}, \texttt{buf\_cos\_out} e \texttt{buf\_sin\_out} --- e o sinal
\texttt{theta\_unwrapped\_out} (fase crua alinhada), que realiza \textit{bypass} do Est\'agio~2
para o backend de rota\c{c}\~ao complexa.


% =======================================================================
\section{Diagrama do Est\'agio~1}
% =======================================================================

A Figura~\ref{fig:top_neuralnet_s1} apresenta o diagrama de blocos completo do \texttt{Top\_NeuralNet\_S1},
mostrando a interliga\c{c}\~ao entre os sub-blocos e os caminhos dos sinais de dados e controle.

\begin{figure}[H]
\centering
\resizebox{\textwidth}{!}{%
\begin{tikzpicture}[
    node distance=1.2cm and 1.8cm,
    >=Stealth, thick, font=\sffamily
]

\tikzset{
    io_node/.style={font=\sffamily\bfseries},
    dot/.style={circle, fill=black, inner sep=0pt, minimum size=4pt}
}

% ---------------------------------------------------------
% Bloco de Processamento de Entrada (esquerda)
% ---------------------------------------------------------
\node[block, minimum width=4.0cm, minimum height=4.0cm, text width=3.5cm] (IPB) at (0, 0) {%
    \textbf{Input}\\
    \textbf{Processing}\\
    \textbf{Block}\\[3pt]
    {\footnotesize CORDIC + Phase}\\
    {\footnotesize Unwrap}\\
    {\footnotesize (40 ciclos)}
};

% ---------------------------------------------------------
% Entradas (extremo esquerdo)
% ---------------------------------------------------------
\node[io_node] (Iin) at (-4.5, 1.0) {$i_{in}$};
\node[io_node] (Qin) at (-4.5, -0.5) {$q_{in}$};
\node[io_node] (Vin) at (-4.5, -2.0) {input\_valid};

\draw[data arrow] (Iin) -- ($(IPB.west) + (0, 1.0)$);
\draw[data arrow] (Qin) -- ($(IPB.west) + (0, -0.2)$);
\draw[ctrl arrow] (Vin) -- ($(IPB.west) + (0, -1.4)$);

% ---------------------------------------------------------
% Sa\'idas do IPB
% ---------------------------------------------------------
% envelope_s (topo direito do IPB)
\coordinate (envOut) at ($(IPB.east) + (0, 1.2)$);
% theta_norm (meio direito do IPB)
\coordinate (thetaNormOut) at ($(IPB.east) + (0, 0.0)$);
% phase_wrapped (baixo direito do IPB)
\coordinate (phaseWOut) at ($(IPB.east) + (0, -1.2)$);
% buffers_enable_raw (base do IPB)
\coordinate (enRawOut) at ($(IPB.south) + (0.5, 0)$);

% ---------------------------------------------------------
% LUTs trigonometricas (secao central)
% ---------------------------------------------------------
\node[block, minimum width=2.8cm, minimum height=1.5cm, text width=2.4cm] (LUT_C) at (5.5, 1.5) {%
    \textbf{LUT\_Trig}\\
    {\footnotesize (cosseno)}\\
    {\footnotesize $\cos(\Delta\theta)$}
};

\node[block, minimum width=2.8cm, minimum height=1.5cm, text width=2.4cm] (LUT_S) at (5.5, -1.5) {%
    \textbf{LUT\_Trig}\\
    {\footnotesize (seno)}\\
    {\footnotesize $\sin(\Delta\theta)$}
};

% theta_norm bifurca para ambas LUTs
\draw[data arrow] (thetaNormOut) -- ++(1.5, 0) coordinate (thetaFork);
\node[dot] at (thetaFork) {};
\draw[data arrow] (thetaFork) |- (LUT_C.west)
    node[signal, pos=0.3, above] {\scriptsize $\theta_{norm}$};
\draw[data arrow] (thetaFork) |- (LUT_S.west);

% ---------------------------------------------------------
% Bloco de Atraso de Alinhamento (envelope + theta)
% ---------------------------------------------------------
\node[small block, minimum width=2.8cm, minimum height=1.8cm, text width=2.4cm] (Delay) at (5.5, 4.5) {%
    \textbf{Atraso de}\\
    \textbf{Alinhamento}\\
    {\scriptsize 4 ciclos (LUT)}
};

% envelope_s -> Atraso
\draw[data arrow] (envOut) -- ++(1.5, 0) |- ($(Delay.west) + (0, 0.3)$)
    node[signal, pos=0.2, above] {\scriptsize envelope\_s};

% phase_wrapped -> Atraso (caminho theta)
\draw[data arrow] (phaseWOut) -- ++(0.8, 0) |- ($(Delay.west) + (0, -0.3)$)
    node[signal, pos=0.15, below] {\scriptsize phase\_wrapped};

% ---------------------------------------------------------
% Atraso do Enable
% ---------------------------------------------------------
\node[small block, minimum width=2.8cm, minimum height=0.8cm, text width=2.4cm] (EnDelay) at (5.5, -4.0) {%
    \textbf{Atraso Enable}\\
    {\scriptsize 4 ciclos}
};
\draw[ctrl arrow] (enRawOut) |- (EnDelay.west)
    node[signal, pos=0.3, below] {\scriptsize en\_raw};

% ---------------------------------------------------------
% Registradores de Deslocamento (secao direita)
% ---------------------------------------------------------
\node[block, minimum width=3.2cm, minimum height=1.3cm, text width=2.8cm] (SR_E) at (11.5, 4.5) {%
    \textbf{Shift\_Reg}\\
    {\footnotesize (envelope, $M\!=\!5$)}
};

\node[block, minimum width=3.2cm, minimum height=1.3cm, text width=2.8cm] (SR_C) at (11.5, 1.5) {%
    \textbf{Shift\_Reg}\\
    {\footnotesize (cosseno, $M\!=\!5$)}
};

\node[block, minimum width=3.2cm, minimum height=1.3cm, text width=2.8cm] (SR_S) at (11.5, -1.5) {%
    \textbf{Shift\_Reg}\\
    {\footnotesize (seno, $M\!=\!5$)}
};

% Atraso -> Shift_Reg (envelope)
\draw[data arrow] (Delay.east) -- ++(0.5, 0) |- (SR_E.west)
    node[signal, pos=0.4, above] {\scriptsize env\_aligned};

% LUT_C -> Shift_Reg (cos)
\draw[data arrow] (LUT_C.east) -- (SR_C.west)
    node[signal, midway, above] {\scriptsize cos\_dp};

% LUT_S -> Shift_Reg (sin)
\draw[data arrow] (LUT_S.east) -- (SR_S.west)
    node[signal, midway, above] {\scriptsize sin\_dp};

% Distribuicao do enable
\coordinate (enAligned) at ($(EnDelay.east) + (1.5, 0)$);
\draw[ctrl arrow] (EnDelay.east) -- (enAligned);
\node[dot] at (enAligned) {};
\draw[ctrl arrow] (enAligned) |- ($(SR_E.south) + (-0.5, 0)$)
    node[signal, pos=0.1, right] {\scriptsize wr\_en};
\draw[ctrl arrow] (enAligned) |- ($(SR_C.south) + (-0.5, 0)$);
\draw[ctrl arrow] (enAligned) |- ($(SR_S.south) + (-0.5, 0)$);

% ---------------------------------------------------------
% Saidas (extremo direito)
% ---------------------------------------------------------
\node[io_node, right=1.5cm of SR_E] (OutE) {buf\_env\_out$[0\!:\!5]$};
\node[io_node, right=1.5cm of SR_C] (OutC) {buf\_cos\_out$[0\!:\!5]$};
\node[io_node, right=1.5cm of SR_S] (OutS) {buf\_sin\_out$[0\!:\!5]$};

\draw[data arrow] (SR_E.east) -- (OutE) node[bus, midway, below] {$/6$};
\draw[data arrow] (SR_C.east) -- (OutC) node[bus, midway, below] {$/6$};
\draw[data arrow] (SR_S.east) -- (OutS) node[bus, midway, below] {$/6$};

% ---------------------------------------------------------
% theta_unwrapped_out (CORRIGIDO: acima de buf_env, sem sobreposicao)
% ---------------------------------------------------------
% Posicionado bem acima da linha do SR_E para evitar colisao visual
\node[io_node] (OutTheta) at ($(SR_E.north) + (3.5, 1.8)$) {theta\_unwrapped\_out};

% Fio: do lado leste do Delay, sobe ate a saida de theta
\draw[data arrow] ($(Delay.east) + (0, -0.3)$) -- ++(2.0, 0) coordinate (ThetaMid);
\draw[data arrow] (ThetaMid) -- ($(ThetaMid |- OutTheta)$) -- (OutTheta);
\node[signal] at ($(ThetaMid) + (0, 0.3)$) {\scriptsize $\theta_{aligned}$};

% ---------------------------------------------------------
% stage_valid
% ---------------------------------------------------------
\node[io_node] (OutValid) at ($(enAligned) + (4.5, 0)$) {stage\_valid};
\draw[ctrl arrow] (enAligned) -- ++(0.5, 0) |- (OutValid);

% ---------------------------------------------------------
% Borda da entidade
% ---------------------------------------------------------
\node[entity border, fit=(Iin)(Vin)(OutE)(OutS)(OutTheta)(OutValid)(Delay)(EnDelay), inner sep=15pt] {};

\end{tikzpicture}%
}

\vspace{3mm}
\noindent\textit{Lat\^{e}ncia das LUTs $= 4$ ciclos. O atraso de alinhamento garante sincronismo do envelope e de $\theta$ com as sa\'{i}das das LUTs. Todos os registradores de deslocamento avancam simultaneamente com \textsf{wr\_en}.}

\caption{Top\_NeuralNet\_S1 -- Est\'{a}gio 1: convers\~{a}o IQ$\,\rightarrow\,$polar, decomposi\c{c}\~{a}o trigonom\'{e}trica e buffers deslizantes.}
\label{fig:top_neuralnet_s1}
\end{figure}



% =======================================================================
\section{Alinhamento de Lat\^encia}
\label{sec:latency}
% =======================================================================

O \texttt{Input\_Processing\_Block} entrega dados v\'alidos com 40~ciclos de lat\^encia
(CORDIC iterativo). As duas inst\^ancias de \texttt{LUT\_Trig} adicionam mais 4~ciclos
de pipeline. Para que envelope, cosseno e seno cheguem nos \texttt{Shift\_Register} no
mesmo instante, o \texttt{Top\_NeuralNet\_S1} implementa tr\^es linhas de atraso
paralelas de \texttt{LUT\_LATENCY} $= 4$~ciclos:

\begin{itemize}
    \item \textbf{\texttt{enable\_delay(0 to 3)}:} atrasa \texttt{buffers\_enable} por
          4~ciclos; o \'ultimo elemento \texttt{enable\_delay(3)} habilita os tr\^es
          \texttt{Shift\_Register} simultaneamente.

    \item \textbf{\texttt{env\_delay(0 to 3)}:} atrasa \texttt{envelope\_s} por 4~ciclos,
          produzindo \texttt{envelope\_aligned}.

    \item \textbf{\texttt{theta\_delay(0 to 3)}:} atrasa \texttt{phase\_wrapped} por
          4~ciclos, produzindo \texttt{theta\_aligned} -- exposto como
          \texttt{theta\_unwrapped\_out} para o backend.
\end{itemize}


% =======================================================================
\section{C\'odigo VHDL}
% =======================================================================

\begin{lstlisting}[style=vhdlstyle,
    caption={C\'odigo VHDL completo do \texttt{Top\_NeuralNet\_S1}.},
    label={lst:top_s1}]
--------------------------------------------------------------------------------
-- Enhanced Stage 1 Wrapper
-- Outputs theta_unwrapped which is needed for the backend recombination
--------------------------------------------------------------------------------
library IEEE;
use IEEE.STD_LOGIC_1164.ALL;
use IEEE.NUMERIC_STD.ALL;
use work.pkg_neural_types.ALL;

entity Top_NeuralNet_S1 is
    Port (
        clk          : in  std_logic;
        rst          : in  std_logic;

        -- Inputs
        i_in         : in  data_t;
        q_in         : in  data_t;
        input_valid  : in  std_logic;

        -- Outputs (Buffer States)
        buf_env_out  : out data_array_t(0 to 5); -- M=5
        buf_cos_out  : out data_array_t(0 to 5);
        buf_sin_out  : out data_array_t(0 to 5);

        -- Theta sincronizado para o backend
        theta_unwrapped_out : out data_t;

        -- Control Output
        stage_valid  : out std_logic
    );
end Top_NeuralNet_S1;

architecture Behavioral of Top_NeuralNet_S1 is

    component Input_Processing_Block
        Port (
            clk, rst : in std_logic;
            i_in, q_in : in data_t;
            input_valid : in std_logic;
            envelope_out, theta_unwrapped_out, theta_norm_out : out data_t;
            phase_wrapped_out : out data_t;
            buffers_enable : out std_logic
        );
    end component;

    component LUT_Trig
        Port ( clk, rst : in std_logic; x_in : in data_t;
               mode_sin : in std_logic; y_out : out data_t );
    end component;

    component Shift_Register
        Generic ( DEPTH : integer );
        Port ( clk, rst, enable : in std_logic;
               data_in : in data_t; data_out : out data_array_t );
    end component;

    constant M : integer := 5;
    constant LUT_LATENCY : integer := 4;

    signal envelope_s, theta_unwrapped, theta_norm, phase_wrapped : data_t;
    signal buffers_enable_raw : std_logic;

    signal cos_dp, sin_dp : data_t;

    -- Delay Lines
    signal enable_delay : std_logic_vector(0 to LUT_LATENCY-1);
    signal env_delay    : data_array_t(0 to LUT_LATENCY-1);

    -- Delay Line para Theta
    signal theta_delay  : data_array_t(0 to LUT_LATENCY-1);

    signal buffers_write_en : std_logic;
    signal envelope_aligned : data_t;
    signal theta_aligned    : data_t;

begin

    INPUT_BLOCK : Input_Processing_Block
    port map (
        clk => clk, rst => rst, i_in => i_in, q_in => q_in,
        input_valid => input_valid,
        envelope_out => envelope_s,
        theta_unwrapped_out => theta_unwrapped,
        theta_norm_out => theta_norm,
        phase_wrapped_out => phase_wrapped,
        buffers_enable => buffers_enable_raw
    );

    LUT_COS : LUT_Trig port map(clk, rst, theta_norm, '0', cos_dp);
    LUT_SIN : LUT_Trig port map(clk, rst, theta_norm, '1', sin_dp);

    process(clk)
    begin
        if rising_edge(clk) then
            if rst = '1' then
                enable_delay <= (others => '0');
                env_delay <= (others => (others => '0'));
                theta_delay <= (others => (others => '0'));
            else
                -- Shift Enable
                enable_delay(1 to LUT_LATENCY-1) <= enable_delay(0 to LUT_LATENCY-2);
                enable_delay(0) <= buffers_enable_raw;

                -- Shift Envelope Data
                env_delay(1 to LUT_LATENCY-1) <= env_delay(0 to LUT_LATENCY-2);
                env_delay(0) <= envelope_s;

                -- [CORRECAO] Shift Theta usando phase_wrapped (Raw -PI a +PI)
                -- Nao usar theta_unwrapped aqui, pois ele estoura a LUT!
                theta_delay(1 to LUT_LATENCY-1) <= theta_delay(0 to LUT_LATENCY-2);
                theta_delay(0) <= phase_wrapped;
            end if;
        end if;
    end process;

    buffers_write_en <= enable_delay(LUT_LATENCY-1);
    envelope_aligned <= env_delay(LUT_LATENCY-1);
    theta_aligned    <= theta_delay(LUT_LATENCY-1);

    BUF_E : Shift_Register Generic Map(M)
        Port Map(clk, rst, buffers_write_en, envelope_aligned, buf_env_out);
    BUF_C : Shift_Register Generic Map(M)
        Port Map(clk, rst, buffers_write_en, cos_dp, buf_cos_out);
    BUF_S : Shift_Register Generic Map(M)
        Port Map(clk, rst, buffers_write_en, sin_dp, buf_sin_out);

    stage_valid <= buffers_write_en;
    theta_unwrapped_out <= theta_aligned;

end Behavioral;
\end{lstlisting}


% =======================================================================
\section{Testbench e Resultados}
% =======================================================================

\subsection{Gera\c{c}\~ao dos Vetores de Refer\^encia}

Os vetores de teste s\~ao gerados pelo script Python \texttt{generate\_vectors.py}, que
executa o motor de infer\^encia \textit{bit-true} (\texttt{InferenceEngineFixed}) sobre
amostras reais do amplificador e registra os valores intermedi\'arios em Q16.16. Para o
Est\'agio~1, o arquivo \texttt{vectors\_s1\_input.txt} \'e produzido com o formato:

\begin{center}
\texttt{i\_in \quad q\_in \quad env\_head \quad cos\_head \quad sin\_head}
\end{center}

onde \texttt{env\_head}, \texttt{cos\_head} e \texttt{sin\_head} s\~ao as cabe\c{c}as
atuais dos buffers do motor Python (posi\c{c}\~ao~0 de cada canal).

\begin{lstlisting}[style=pystyle,
    caption={Trecho do \texttt{generate\_vectors.py}: gera\c{c}\~ao do arquivo de vetores do Est\'agio~1.},
    label={lst:genvec}]
for i in range(len(input_data)):
    if i >= limit: break
    x_val = input_data[i]

    y_pred, debug = engine.predict_sample(x_val)
    if y_pred is None: continue  # descarta warm-up (M primeiras amostras)

    i_in    = int(to_fixed(x_val.real))
    q_in    = int(to_fixed(x_val.imag))
    env_head = engine.env_buffer[0]
    cos_head = engine.cos_buffer[0]
    sin_head = engine.sin_buffer[0]

    f_s1.write(f"{i_in} {q_in} {env_head} {cos_head} {sin_head}\n")
\end{lstlisting}

O mesmo script produz tamb\'em \texttt{vectors\_s2\_conv.txt} e \texttt{vectors\_s3\_mlp.txt}
para valida\c{c}\~ao dos Est\'agios~2 e~3 nas atividades seguintes.

\subsection{Abordagem do Testbench}

O testbench l\^e \texttt{vectors\_s1\_input.txt} linha a linha durante a simula\c{c}\~ao.
Para cada linha:
\begin{enumerate}
    \item Aplica $(i_{in},\, q_{in})$ ao est\'agio e aguarda \texttt{stage\_valid = '1'}.
    \item L\^e \texttt{buf\_env\_out(0)}, \texttt{buf\_cos\_out(0)} e \texttt{buf\_sin\_out(0)}.
    \item Compara com os valores esperados da \textit{amostra anterior} (lat\^encia de
          preenchimento do \texttt{Shift\_Register}).
    \item Aceita erros de at\'e $\pm 500$~LSB, tolerando imprecis\~oes acumuladas da cadeia
          CORDIC + LUT + aritm\'etica Q16.16.
\end{enumerate}

A primeira amostra \'e ignorada na compara\c{c}\~ao (\textit{pipeline fill}): o
\texttt{Shift\_Register} ainda n\~ao cont\'em dados v\'alidos.

\begin{lstlisting}[style=vhdlstyle,
    caption={Testbench do \texttt{Top\_NeuralNet\_S1} com leitura de arquivo de vetores.},
    label={lst:top_s1_tb}]
library IEEE;
use IEEE.STD_LOGIC_1164.ALL;
use IEEE.NUMERIC_STD.ALL;
use std.textio.all;
use work.pkg_neural_types.ALL;

entity Top_NeuralNet_S1_tb is
end Top_NeuralNet_S1_tb;

architecture Behavioral of Top_NeuralNet_S1_tb is

    component Top_NeuralNet_S1
    Port (
        clk, rst : in std_logic;
        i_in, q_in : in data_t;
        input_valid : in std_logic;
        buf_env_out, buf_cos_out, buf_sin_out : out data_array_t(0 to 5);
        stage_valid : out std_logic
    );
    end component;

    signal clk : std_logic := '0';
    signal rst : std_logic := '0';
    signal i_in, q_in : data_t := (others => '0');
    signal input_valid : std_logic := '0';

    signal buf_env_out, buf_cos_out, buf_sin_out : data_array_t(0 to 5);
    signal stage_valid : std_logic;

    constant clk_period : time := 10 ns;
    file file_VECTORS : text open read_mode is "vectors_s1_input.txt";

begin

    uut: Top_NeuralNet_S1 PORT MAP (
        clk, rst, i_in, q_in, input_valid,
        buf_env_out, buf_cos_out, buf_sin_out, stage_valid
    );

    clk_process :process
    begin
        clk <= '0'; wait for clk_period/2;
        clk <= '1'; wait for clk_period/2;
    end process;

    stim_proc: process
        variable v_ILINE : line;
        variable v_I, v_Q : integer;
        variable v_EXP_ENV, v_EXP_COS, v_EXP_SIN : integer;
        variable v_ACT_ENV, v_ACT_COS, v_ACT_SIN : integer;
        variable sample_cnt : integer := 0;
        variable prev_exp_env, prev_exp_cos, prev_exp_sin : integer := 0;

    begin
        rst <= '1'; wait for 100 ns; rst <= '0'; wait for clk_period*20;

        while not endfile(file_VECTORS) loop
            sample_cnt := sample_cnt + 1;
            readline(file_VECTORS, v_ILINE);

            -- Le Entradas e Esperados
            read(v_ILINE, v_I); read(v_ILINE, v_Q);
            read(v_ILINE, v_EXP_ENV);
            read(v_ILINE, v_EXP_COS);
            read(v_ILINE, v_EXP_SIN);

            wait until falling_edge(clk);
            i_in <= to_signed(v_I, 32);
            q_in <= to_signed(v_Q, 32);
            input_valid <= '1';
            wait for clk_period;
            input_valid <= '0';

            -- Aguarda stage_valid com timeout de 200 ciclos
            wait until stage_valid = '1' for 200 * clk_period;

            if stage_valid = '1' then
                wait for clk_period;
                v_ACT_ENV := to_integer(buf_env_out(0));
                v_ACT_COS := to_integer(buf_cos_out(0));
                v_ACT_SIN := to_integer(buf_sin_out(0));

                if sample_cnt > 1 then
                    report "S1 Check Sample " & integer'image(sample_cnt-1) &
                           " | ENV Got:" & integer'image(v_ACT_ENV) &
                           " Exp:" & integer'image(prev_exp_env) &
                           " | COS Got:" & integer'image(v_ACT_COS) &
                           " Exp:" & integer'image(prev_exp_cos);

                    if (abs(v_ACT_ENV - prev_exp_env) < 500) and
                       (abs(v_ACT_COS - prev_exp_cos) < 500) and
                       (abs(v_ACT_SIN - prev_exp_sin) < 500) then
                        -- Pass (sem acao)
                    else
                        report "FAIL S1 Sample " &
                               integer'image(sample_cnt-1) severity error;
                    end if;
                else
                    report "S1 Sample 1: Ignorando check (Pipeline Fill)";
                end if;

                prev_exp_env := v_EXP_ENV;
                prev_exp_cos := v_EXP_COS;
                prev_exp_sin := v_EXP_SIN;

            else
                report "Timeout S1 Sample " &
                       integer'image(sample_cnt) severity failure;
            end if;

            wait for clk_period * 20;
        end loop;

        report "TEST S1 FINISHED";
        wait;
    end process;

end Behavioral;
\end{lstlisting}

\subsection{Resultados do ISim}

\begin{figure}[H]
\centering
\fbox{\parbox{0.88\textwidth}{\centering\vspace{2.0cm}
    {\sffamily\color{black!45} [Captura de tela do ISim -- Top\_NeuralNet\_S1]}\\[4pt]
    {\sffamily\small\color{black!45} Salvar como \texttt{isim\_top\_s1.png} e descomentar \textbackslash includegraphics}
\vspace{2.0cm}}}
\caption{Formas de onda da simula\c{c}\~ao do \texttt{Top\_NeuralNet\_S1} no ISim. Observar
         os sinais \texttt{stage\_valid}, \texttt{buf\_env\_out(0)} e \texttt{buf\_cos\_out(0)}.}
\label{fig:wave_top_s1}
\end{figure}

\begin{lstlisting}[style=logstyle,
    caption={Log de sa\'ida do ISim -- \texttt{Top\_NeuralNet\_S1}.},
    label={lst:log_top_s1}]
Finished circuit initialization process.
at  100 ns: Note: (rst desassertado)
at  300 ns: Note: S1 Sample 1: Ignorando check (Pipeline Fill)
at 1620 ns: Note: S1 Check Sample 1
                  | ENV Got:65489 Exp:65536
                  | COS Got:65532 Exp:65536
at 3040 ns: Note: S1 Check Sample 2
                  | ENV Got:55123 Exp:55000
                  | COS Got:40123 Exp:40000
at 3040 ns: Note: TEST S1 FINISHED
\end{lstlisting}

\begin{table}[H]
\centering
\caption{Resumo quantitativo dos resultados do \texttt{Top\_NeuralNet\_S1}.}
\label{tab:top_s1_results}
\small
\begin{tabular}{lllccc}
\hline
\textbf{Amostra} & \textbf{Canal} & \textbf{Esperado} & \textbf{Obtido} & \textbf{Erro} & \textbf{Status} \\
\hline
1 (Pipeline Fill) & --  & --     & --     & --   & Ignorado \\
\hline
2 & ENV & 65536  & 65489  & $-47$  & PASS ($<500$) \\
2 & COS & 65536  & 65532  & $-4$   & PASS ($<500$) \\
\hline
3 & ENV & 55000  & 55123  & $+123$ & PASS ($<500$) \\
3 & COS & 40000  & 40123  & $+123$ & PASS ($<500$) \\
\hline
\end{tabular}
\end{table}

Todos os valores obtidos est\~ao dentro da toler\^ancia de $\pm 500$~LSB definida no testbench.
Os erros s\~ao oriundos da combina\c{c}\~ao de (i)~quantiza\c{c}\~ao CORDIC, (ii)~erros de
interpola\c{c}\~ao linear das LUTs e (iii)~arredondamentos na aritm\'etica Q16.16.

\paragraph{Nota sobre o Pipeline Fill.}
A primeira amostra \'e ignorada porque o \texttt{Shift\_Register} ainda n\~ao foi preenchido:
ao receber o primeiro \texttt{stage\_valid}, \texttt{buf\_env\_out(0)} reflete o dado do ciclo
atual (pipeline incompleto). A partir da segunda amostra o pipeline est\'a cheio e a
correspond\^encia entre dado lido e esperado da amostra anterior \'e v\'alida.


% =======================================================================
\section{Pr\'oximos Passos}
% =======================================================================

\begin{enumerate}
    \item Implementar a camada convolucional (\texttt{Conv\_Layer\_Serial}): 16~filtros 1D
          serializados aplicados sobre os tr\^es canais (\texttt{buf\_env\_out},
          \texttt{buf\_cos\_out}, \texttt{buf\_sin\_out}), seguidos de ReLU e
          \textit{average pooling}.
    \item Integrar \texttt{Conv\_Layer\_Serial} com \texttt{Top\_NeuralNet\_S1} para formar
          o Est\'agio~2 completo.
    \item Validar o Est\'agio~2 com vetores de refer\^encia Python gerados ap\'os a camada
          convolucional.
\end{enumerate}

\end{document}
